\documentclass[10pt]{article}

\usepackage[utf8]{inputenc}

\usepackage[T1]{fontenc}

\usepackage{amsmath}

\usepackage{amsfonts}

\usepackage{amssymb}

\usepackage[version=4]{mhchem}

\usepackage{stmaryrd}



\begin{document}

переноса п т. Д., а также исследовать особенности простейших возбуждений системы (в общем случае при $\theta \neq 0)$, их энергию, затухание и т. д.



Указанные средние расснитываются до конца только в исключительных случаях: для идеальных систем и для нек-рых специальных моделөй. Эти расчеты в дальнейших исследованиях могут служить нулевым приближением. Наиболее часто рассматриваемыми моделями неидеальных статистич. систем являются системы с прямым взаимодействием частиц друг с другом (взаимодействием конечного радиуса, кулоновским взаимодействием и др.), с взаимодействием частцц с полем типа фотонного (в твердом теле описывающего тепловое движение кристаллич. решетки), дискретные системы типа Изинга и гейзенберговского магнетика с взаимодействием узлов конечного радиуса действия, а такæе сочетания взаимодействий подобных типов. В представлении вторичного квантования гамильтониан $H=$ $=H_{0}+H_{1}$ в этих случаях выразается чөрез квадратичные комбинации операторов рождения И уничтожения (в части $H_{0}$ без взаимодействия), четвертную форму (если $H_{1}$ включает прямое взаимодействие частиц), тройную форму тина используемой в гвантовой электродинамике (әлектранфотонное взаимодействие в $H_{1}$ ) и т. Д.



Приближенные методы расчета указанных средних в большинстве случаев основываются на добавлении поправок к результатам, полученцым для случая $H=H_{0}$ (если нулевое приближение в физич. отношении действительно является таковым), имеющих вид явного или модифицированного разложения по степеням параметра, определяющего интенсивность взаимодействия, включенного в гамильтониан $H_{1}$. При сопоставлении рассматриваемой формальной модели с реальными системами в целом ряде случаев, имеющих прикладной интерес, параметр взаимодействия, по к-рому производится "разложение», не оказывается малым. Такие трудности существуют в проблеме электронного газа в металлах, в теории неидеального бозе-газа, в теории. жидкостеӥ, при расспотрении ситуаций в области фазовых переходов или вблизя критич. точек и т. Д.



Если отвлечься от этих проблем постановочного характера и считать, тто малый параметр для каких-то частных случаев все же ияеется, то при развитии теории возмущениї по нему появляются специфические для многотельных систем трудности, формально проявляющиеся в появлении расходимостей в членах, учитывающих многочастичные корреляции. Их Іоявление связано с тем, что простой ряд по целым степеням используемого «малого» параметра, начиная c нек-рой степени, уже не отражает действительной зависимости исследуемой характеристики системы от этого параметра. Эти трудности, как чисто математические, в принципе преодолимы. Для выявления этих «неаналитически» зависящих от малого параметра поправок разработаны методы, в основе своей связанные с учетом определенных классов наиболее существенных для каждого конкретного случая формального ряда теории возмущений по вараметру взаимодействия.



В задачах статистич. механики классич. систем с прямым взаимодействием частиц исследование в большинстве случаев строится на основе развития метода Боголюбова [1] (см. Боголюбова цепочка уравнений) или метода Майера [2]. В первом случае, связанном с исследованием цепочки интегро-дифференциальных уравнений для одно-, двух- и т. Д. частичных функций распределения, основой аппроксимационной процедуры является обрыв этой цепочки. При әтом входящая в интегральную часть последнего из рассматриваемых уравнений цепочки функция распределения более высокого ранга выражается с помощью какой-либо комбинации функций распределения низшего порядка. Процедура такого расцепления исходит из анализа физич. особенностей системы и характерного для данной ситуации малого параметра. Так, в системах с короткодействующими силами взаимодействия - это куб отношения радиуса взаимодействия к среднему расстоянию мөжду частицапи (основа т. н. вириального разложения, или разложения по обратным степеням удельного объеда $v=V / N$ ); в системах с кулоновским взаимодействием такой параметр связан с отношением средней энергии кулоновского взаимодействия частиц к средней их кинетиף. энергии и т. д. После расцепления цепочки математич. проблема сводится к решению системы нелинейных в своей интегральной части уравнөний (или одного нелинейного интегрального уравнения) относительно функций распределения, подчиненных условиям нормировки и ослабления корреляций (распадение многочастичных функций при «раздвигании» пространственных их аргументов на произведение функций более низкой частичности), играющих роль своеобразных граничңых условий.



М е т о д М а й е р а (см. [2]) для систем с короткодействующим потенциалом взаимодействия



\[

\Phi_{i j}=\Phi\left(\left|\boldsymbol{r}_{i}-\boldsymbol{r}_{j}\right|\right)

\]



включаюццм бесконечное отталкивание на малых расстояниях, исходит из представления классич. интеграла састояний $Z$ в виде



\[

\frac{Z}{Z_{\text {ид }}}=\frac{1}{V^{N}} \int_{(V)} d \boldsymbol{r}_{1} \cdots d \boldsymbol{r}_{N} \prod_{1 \leqslant i<i \leqslant N}\left(1+f_{i j}\right),

\]



где фу ук ци я $\mathrm{M}$ а ї е р a $f_{i j}=\exp \left(-\Phi_{i j} / \theta\right)-1$, и получения для удельной свободної энергии



\[

F / N=-\theta \ln Z^{1 / N}

\]



и других характеристик системы вириальных разложений с коэффициентами, выражающимися через интегралы от произведений все возрастающего числа функций $f_{i j}$, связанных попарно одним пз двух своих аргументав $\boldsymbol{r}_{j}$ или $\boldsymbol{r}_{i}$ (для реальных моделей потенциала $\Phi_{i j}$ эти интегралы вычисляются с помощью численных методов). Математич. проблемы, возникающие в связи с использованием вириального разложения, - это проблемы разработки методов суммирования этих рядов (хотя бы частичного) с целью подойти к области фазового перехода газ - жидкость или критич. области, а также общие проблемы сходимости самого разложения и т. Д.



В задачах статистич. механики неидеальных квантовых систем, связанной с представлением операторов динамич. величин через квантовые операторы рождения или уничтожения, эфффективными оказались методы, развитые первоначально в квантовой теории поля. Наиболее распространенными являготся бестемпературная многовременная техника функций Грина причинного типа, безвременная температурная техника и метод двухвременных температурных функций Грина. В первых двух подходах (см. [3]) формальный ряд теории возмущений по интенсивности взаимодействия частиц друг с другом или каким-либо полем в известном смысле аналогичен соответствующим разложениям в квантовой теории поля (в чисто температурной технике роль «времени» играет мнимая обратная температура). 'Поэтому значительное развитие в этих формализмах получили диаграммные представления. ряда теории возмущений, сводящиеся к рассмотрению вмосто исходных частиц, и взаимодействия - квазичастиц с перенормированной энергией и конечным затуханием, перенормированного әффективного взаимодействия (вершинных частей) и т. Д. Система интегральных уравнений для функции Грина квазичастиц Ir





\end{document}